The DMA controller (\peripheral{DMAx}) is a configurable multi-channel single-shot direct-memory-access engine. It moves words from a source to a destination over the shared multi-core arbiter as a stream of single-word transactions --- either flat-out under software command (memory-to-memory) or paced one word per peripheral data-ready event --- with per-channel source and destination auto-increment, a two-level priority with word-granular round-robin between channels, an optional \mbox{CRC16-CDMA2000} ride-along accumulated over the transferred data, and a hardware read-side-effect guard. The channel count is a build-time generic (\{2, 4\}); the register map is always the four-channel superset, so on a two-channel build the upper channel blocks and their status bits read 0. The whole transfer engine --- the master-port state machine, the per-channel working counters, the round-robin picker, the CRC datapath, the pacing edge-detectors, the sticky event flags, and the two interrupt combiners --- rides the free-running fabric clock (MCLK), while the register file rides the gated bus clock. It has \emph{no pins}. The main features are listed below:

\begin{itemize}
	\item Two or four independent single-shot channels, each programmed with a source byte address, a destination byte address, a word length, and a configuration word (\register{DMAxCnSRC}, \register{DMAxCnDST}, \register{DMAxCnLEN}, \register{DMAxCnCFG})
	\item Memory-to-memory transfers under software command, or peripheral-paced transfers that move one word per data-ready event from the on-chip \peripheral{UARTx}, \peripheral{QSPIx}, or \peripheral{NFCx} receivers (\register{DMACnTRIG})
	\item Per-channel source and destination auto-increment (\register{DMACnSINC}, \register{DMACnDINC}): a fixed pointer drains or fills a peripheral data register, an incrementing pointer walks a buffer
	\item A two-level priority class (\register{DMACnPRIO}) with round-robin scheduling within a class, servicing one word per channel visit so no channel in a class starves
	\item An optional \mbox{CRC16-CDMA2000} ride-along (\register{DMACnCRCEN}) folding each transferred word into the shared \register{DMAxCRC} accumulator, for an integrity check that costs no extra bus traffic
	\item A same-cycle \register{DMABUSY} status and sticky write-1-to-clear per-channel done and error flags (\register{DMADONE}, \register{DMAERR}), with two lean interrupts: combined done (vector 118) and error (vector 119)
	\item A hardware read-side-effect guard that refuses engine reads of the mutex bank or the interrupt-router claim word, plus an at-launch check for a zero length, a misaligned or out-of-window pointer, or a pointer into the private-TCM hole
\end{itemize}

\subsection{Programming and Launching a Transfer}

A channel is programmed by writing its source and destination byte addresses, its word length, and its configuration word, then launched by a byte-lane-0 write to \register{DMAxCR} that sets the channel's \register{DMAGO} command bit. \register{DMAGO} is a self-clearing command (it reads 0); the write captures the channel's programmed \{source, destination, length, configuration\} into the engine's working registers and starts it. The launch is suppressed --- the content is still captured but no transfer begins --- when the master enable \register{DMAEN} is clear or when that channel is already busy, so a stray second \register{DMAGO} cannot reload a channel mid-transfer. The rule is one outstanding \register{DMAGO} per channel: poll \register{DMABUSY} clear (or wait for the channel's \register{DMADONE}) before re-arming it.

\register{DMABUSY} in \register{DMAxSR} covers the launch instant with no blind window: it reads 1 in the same cycle as the launching \register{DMAGO} write (any channel active \emph{or} a launch pending), and stays high through real engine activity, so firmware may write \register{DMAGO} and immediately spin on \register{DMABUSY} with no chance of falling through. \register{DMAACTIVECH} reports the index of the channel currently being serviced (0 when idle). Source and destination addresses read back the programmed values; \register{DMAxCnLEN} reads the current remaining word count, decrementing as the engine runs and reading 0 before the first launch --- so after an abort the programmed pointers plus the remaining length show exactly where the transfer stopped.

Each transferred word is a read transaction followed by a write transaction, issued as two separate single-word arbiter transactions with an observed idle gap between them. The engine never holds the arbiter grant and never bursts, so it can never pin the shared bus; it participates as one additional fair master and cannot starve the harts. A direct consequence, which firmware must respect, is that a DMA word copy is \emph{not} atomic: another master may write the source or destination word between the read and the write of the same word. While a channel is active, ownership of its source and destination buffers is a software contract --- do not let another agent race the region the DMA is walking.

\subsection{Peripheral-Paced Transfers}

Setting \register{DMACnTRIG} to a nonzero value paces the channel off a peripheral's receive data-ready event instead of running flat out: value 1 paces off the \peripheral{UARTx} receive-complete event, 2 off the \peripheral{QSPIx} receive-full event, and 3 off the \peripheral{NFCx} payload-ready event. Each qualifying event authorizes exactly one word (one read/write pair), with the source pointer normally held fixed (\register{DMACnSINC}$=0$) so successive events drain the same peripheral data register into a growing buffer. Pacing requires the source peripheral's own receive interrupt-enable bit to be set (the engine taps the peripheral's enabled interrupt level, not a raw flag); \emph{routing} that peripheral's interrupt to a hart through the interrupt router is independent and optional --- an enabled-but-unrouted source that only feeds the DMA is a normal and benign configuration.

The three sources differ in how their ready flag clears, and the engine handles each:

\begin{itemize}
	\item \peripheral{UARTx} receive is \emph{read-to-clear}: the engine's own read of the receive register acknowledges the byte and re-arms the flag, so no extra transaction is needed. Each word carries one received byte in the low lane.
	\item \peripheral{QSPIx} receive-full is \emph{side-effect-free} on read, so after draining the word the engine issues one extra write transaction that writes-1-to-clear the receive-full bit in the \peripheral{QSPIx} status register (status slot 5, at offset \texttt{0x14} within the peripheral) to re-arm the source. This is the only case that costs an extra bus transaction per word.
	\item \peripheral{NFCx} payload-ready is \emph{frame-granular}: it fires once per received reader frame, not once per byte, and a payload read auto-advances the peripheral's own buffer index rather than clearing the flag. So one frame event launches the channel's full programmed-length burst (the engine drains the whole length from the fixed payload register), then a single write-1-to-clear acknowledges the frame. Firmware sizes the length to the frame from the \peripheral{NFCx} receive-status register.
\end{itemize}

\subsection{The CRC Ride-Along}

Setting \register{DMACnCRCEN} folds every word that channel transfers into the shared \register{DMAxCRC} accumulator (\mbox{CRC16-CDMA2000}, polynomial \texttt{0xC857} --- the same engine as the SYSTEM peripheral's CRC), feeding the four bytes of each word low byte first. There is one accumulator shared across all channels and no auto-seed: the firmware sequence is seed-then-go. Write \register{DMAxCRC} to the \texttt{0xFFFF} seed, launch the channel, wait for its \register{DMADONE}, then read \register{DMAxCRC} for the result (no input or output reflection, no final XOR). Because the accumulator is shared and the engine interleaves channels word-by-word, a meaningful CRC requires exactly one \register{DMACnCRCEN} channel per measurement session --- two enabled at once accumulate a CRC over the interleaved streams. One further caveat: if a \register{DMACnCRCEN} channel is aborted mid-transfer, its accumulator is left partway through and is not a meaningful check value.

\subsection{Errors, Aborts, and the Read-Side-Effect Guard}

The engine will not issue a read that would have a side effect. Before it would drive a read transaction it compares its own generated word address against a small deny table: the whole mutex-bank window (byte addresses \texttt{0x6000}--\texttt{0x60FF}, where a read atomically claims a mutex) and the interrupt-router claim word (byte address \texttt{0x7800}, where a read atomically claims an interrupt) are denied. A denied read is suppressed (no transaction is issued), the channel's \register{DMAERR} flag sets, the channel aborts, and the error interrupt fires. The router's pending, in-service, and enable registers are side-effect-free and remain legitimately DMA-readable; only the claim word is denied. The guard covers reads only. A DMA \emph{write} into the mutex or router pages would also have a side effect (a mutex release, or an interrupt-complete), and those are \emph{not} guarded in hardware this phase --- keeping a destination pointer out of the mutex and router pages is a software contract.

\register{DMAERR} is also set, with no transfer issued, when a channel is launched with a zero length, with a source or destination that is not word-aligned, with a pointer outside the reachable shared window (at or above the extended-flash base), or with a pointer into the private tile-TCM hole (\texttt{0x8000}--\texttt{0xBFFF}, which is not an arbiter slave). These are checked at the launching \register{DMAGO} and reject the transfer outright rather than silently masking the pointer. There is no bus-fault input from the arbiter, so a bad-but-in-window pointer copies zeros or garbage exactly as a stray CPU load would --- only the out-of-window, misaligned, TCM-hole, zero-length, and deny cases are caught.

A byte-lane-0 write to \register{DMAxCR} that sets a channel's \register{DMAABORT} command bit requests an orderly stop. If that channel owns an in-flight arbiter transaction the engine lets that single transaction complete normally --- the handshake is never truncated --- then drops the channel's busy \emph{without} setting \register{DMADONE} or \register{DMAERR} (an abort is neither a completion nor an error). \register{DMAABORT} is self-clearing (reads 0), and the programmed pointers plus the remaining \register{DMAxCnLEN} show where the transfer stopped.

Both interrupts are combinational and never latched: the done interrupt (vector 118) is the OR of the \register{DMADONE} flags gated by \register{DMADONEIE}, and the error interrupt (vector 119) is the OR of the \register{DMAERR} flags gated by \register{DMAERRIE}. The service routine reads \register{DMAxSR}, discovers which channels' \register{DMADONE} or \register{DMAERR} bits are set, and clears them write-1-to-clear (a set arriving the same cycle as a clear survives). A separate note applies to the shared NPU staging RAM: while an NPU computation is running the NPU owns that RAM, so firmware must not aim a DMA source or destination into the NPU staging window until the NPU reports idle --- a software contract, not a hardware guard.

\subsection{Shared Access on \AsicNameForUserGuide}

In configurations that include it, \peripheral{DMA0} lives in the shared peripheral window behind the multi-core arbiter (see Section \ref{s:multicore}) in the mutex-bank page, sub-slot 8 at \texttt{0x6800}, so any hart can program and launch it. Because the engine clocks off the free-running fabric clock, its progress is immune to clock reconfiguration, and software must \emph{not} apply the ``write \register{\register{SYSCLKCR}} to 0 first'' rule that the SMCLK peripherals need.

\peripheral{DMA0} is unique among the shared-window peripherals in also being an arbiter \emph{master}: its transfer engine is an additional bus master alongside the harts, so enabling the block widens the shared fabric by one master. It participates in the same fair round-robin as the harts, issuing single-word transactions with an idle gap between each, so it is starvation-safe by the arbiter's own fairness and cannot lock the bus.

Its two interrupt sources occupy vectors 118 (combined channels-done) and 119 (error); the done source takes the lower identifier, and both are delivered through the interrupt router like every other shared-peripheral source. These sources sit above the external-interrupt (meip) delivery vector, which stays fixed at vector 85, above the I3C, NFC, \peripheral{GPIOx} port-4/port-5, RTC, PWM, and 1-Wire sources at vectors 86 through 117.
